\section{Extended experimental details}
\label{app:details}

\subsection{Training convergence}

The U-Net with ResNet-34 encoder converges around epoch 60 of 80. We
monitor validation Dice and select the checkpoint with the highest value.
Typical training loss (Dice + BCE hybrid) decreases from $\sim$0.65 at
epoch 1 to $\sim$0.18 at convergence. Validation Dice reaches 0.77 by
epoch 30 and improves by less than 0.01 over the remaining epochs.

Gradient clipping at norm 1.0 is active in roughly 15\% of batches during
the first 10 epochs and less than 2\% after epoch 30; training instability
is confined to early optimization.

\subsection{Vessel-aware patch sampling}

Class imbalance in retinal vessel segmentation is severe: vessels occupy
10--15\% of pixels in DRIVE images. Uniform random patch sampling would
frequently produce patches with no vessels, wasting training signal. Our
vessel-aware sampling constrains 80\% of patches to contain at least 16
vessel pixels (approximately 0.025\% of a 256$\times$256 patch). The
remaining 20\% are uniformly sampled to maintain exposure to background
regions.

This sampling strategy improves sensitivity by 3--5\% compared to uniform
sampling, with negligible effect on specificity. The 16-pixel minimum was
chosen empirically; lower thresholds (4, 8) produced similar results,
while higher thresholds (32, 64) biased training toward thick vessels and
reduced thin-vessel sensitivity.

\subsection{MC Dropout placement}

We apply Dropout2d (spatial dropout) after the decoder and before the
final 1$\times$1 convolution. Alternative placements were tested:

\begin{itemize}[leftmargin=*]
  \item \textbf{After each encoder block:} Produces higher uncertainty
  but lower Unc-AUROC (0.72). Dropout perturbations propagate
  through the entire decoder, producing spatially correlated noise that
  does not localize to error-prone regions.

  \item \textbf{After each decoder block:} Similar Unc-AUROC to our default
  (0.76 vs.\ 0.77) but slower inference due to more dropout operations
  per pass.

  \item \textbf{After decoder only (our choice):} Best trade-off between
  Unc-AUROC and inference speed. A single dropout layer perturbs the
  high-resolution feature map just before the segmentation head, producing
  uncertainty that reflects representation-level disagreement without the
  spatial smoothing caused by decoder upsampling.
\end{itemize}

\subsection{TTA inverse transform verification}

Geometric TTA requires exact inverse transforms. We verify that
$\tau_k^{-1}(\tau_k(\mathbf{y})) = \mathbf{y}$ for every augmentation and
every test image. All 6 augmentations (identity, hflip, vflip, rot90,
rot180, rot270) are exact: they permute pixel indices without
interpolation, so the inverse is also an index permutation. This is not
the case for arbitrary rotations or elastic deformations, which require
interpolation and introduce small reconstruction errors.

\section{Additional results}
\label{app:additional_results}

\subsection{Per-image analysis}

\begin{table}[H]
\centering
\caption{Per-image Dice and Unc-AUROC for 5 representative DRIVE test
images, ranked by difficulty. Harder images (lower Dice) yield higher
Unc-AUROC, consistent with the cross-dataset finding: uncertainty quality
\emph{improves} when the model makes more errors. Best Unc-AUROC in bold.}
\label{tab:per_image}
\scriptsize
\setlength{\tabcolsep}{3.5pt}
\renewcommand{\arraystretch}{1.03}
\begin{center}
\begin{tabular}{l c c c c}
\toprule
\textbf{Image} & \textbf{Dice} & \textbf{Unc-AUROC (MC)} &
\textbf{Unc-AUROC (TTA)} & \textbf{Uncertainty (mean)} \\
\midrule
Hardest & 0.721 & \textbf{0.812} & \textbf{0.913} & 0.019 \\
Q1      & 0.748 & 0.789 & 0.895 & 0.015 \\
Median  & 0.765 & 0.770 & 0.878 & 0.012 \\
Q3      & 0.782 & 0.752 & 0.866 & 0.010 \\
Easiest & 0.809 & 0.718 & 0.842 & 0.007 \\
\bottomrule
\end{tabular}
\end{center}
\end{table}

Table~\ref{tab:per_image} shows a clear inverse relationship between
image difficulty and uncertainty quality. On the hardest image (Dice 0.721),
TTA Unc-AUROC reaches 0.913. On the easiest image (Dice 0.809), Unc-AUROC
drops to 0.842. This pattern is beneficial: the images where deferral
matters most are exactly those where the uncertainty signal is strongest.
Mean uncertainty also correlates with difficulty (0.019 vs.\ 0.007) and
could serve as a simple image-level difficulty estimator for triage.

\subsection{Failure mode analysis}

We categorize segmentation errors into three types and examine how each
uncertainty method handles them.

\paragraph{Boundary errors.} The most common error type. Vessel boundaries
are annotated at single-pixel precision, but the model's soft predictions
span 2--3 pixels across boundaries. TTA flags these regions with high
variance (boundary pixels often change class under rotation or flip);
MC Dropout dilutes the signal by assigning moderate uncertainty to
neighboring non-boundary pixels as well.

\paragraph{Thin vessel misses.} The model sometimes fails to detect vessels
narrower than 2 pixels, particularly in the periphery. Both MC Dropout and
TTA elevate uncertainty in these regions, but TTA is more consistent:
under rotation, a thin vertical vessel aligns differently with horizontal
convolutional filters, producing larger disagreement.

\paragraph{False positives on artifacts.} Occasional bright artifacts
(reflections, drusen-like features) are misclassified as vessels. These
are the most dangerous errors for clinical use because the model tends to
be confident (high $\hat{p}$, low $u$). Neither uncertainty method reliably
catches these. Confidence-aware deferral helps slightly---some artifact
pixels have moderate uncertainty that pushes $s = u \cdot (1-c)$ above
threshold---but this failure mode remains the primary limitation of
purely uncertainty-based deferral.

\subsection{Deferral map visualization}

\begin{figure}[H]
\centering
\includegraphics[width=\textwidth,height=0.38\textheight,keepaspectratio]{deferral_map_easy.png}\\[0.2em]
\includegraphics[width=\textwidth,height=0.38\textheight,keepaspectratio]{deferral_map_hard.png}
\caption{Deferral maps comparing global and confidence-aware strategies.
Confidence-aware deferral adapts to image difficulty and concentrates on
pixels near the decision boundary.}
\label{fig:deferral_maps}
\end{figure}

\subsection{Uncertainty separability}

The Unc-AUROC gap between TTA (0.881) and MC Dropout (0.722) reflects how
cleanly each method's uncertainty separates correct from incorrect pixels.
For MC Dropout, the two distributions overlap heavily: no threshold cleanly
partitions them. For TTA, there exists a threshold range where most
incorrect pixels lie above and most correct pixels lie below, enabling
efficient deferral. This separation is what the main-text risk-coverage
curves (Figure~\ref{fig:risk_coverage}) and decile plot
(Figure~\ref{fig:unc_vs_error}) reflect.

\subsection{Runtime breakdown}

\begin{table}[H]
\centering
\caption{Per-image runtime breakdown on DRIVE (mean over 20 images).
Forward passes dominate total cost; aggregation and deferral scoring are
negligible. Best (fastest) per row in bold.}
\label{tab:runtime}
\scriptsize
\setlength{\tabcolsep}{3.5pt}
\renewcommand{\arraystretch}{1.03}
\begin{center}
\begin{tabular}{l c c c c}
\toprule
\textbf{Component} & \textbf{MC Drop.\ ($T{=}30$)} & \textbf{TTA ($K{=}6$)} &
\textbf{Ensemble ($N{=}5$)} & \textbf{Determ.} \\
\midrule
Preprocessing     & 0.005\,s & 0.005\,s & 0.005\,s & 0.005\,s \\
Forward passes    & 0.780\,s & 0.110\,s & 3.850\,s & \textbf{0.065\,s} \\
Aggregation       & 0.030\,s & 0.012\,s & 0.040\,s & ---      \\
Deferral scoring  & 0.005\,s & 0.003\,s & 0.005\,s & ---      \\
\midrule
\textbf{Total}    & 0.820\,s & 0.130\,s & 3.900\,s & \textbf{0.070\,s} \\
\bottomrule
\end{tabular}
\end{center}
\end{table}

\subsection{Statistical significance}

\begin{table}[H]
\centering
\caption{Paired $t$-test $p$-values for pairwise method comparisons on
the DRIVE test set. The Dice gap between MC Dropout and TTA is not
significant ($p{=}0.142$); the Unc-AUROC gap is highly significant
($p{<}0.001$). ** denotes $p{<}0.01$.}
\label{tab:significance}
\scriptsize
\setlength{\tabcolsep}{3.5pt}
\renewcommand{\arraystretch}{1.03}
\begin{center}
\begin{tabular}{l l c}
\toprule
\textbf{Comparison} & \textbf{Metric} & $p$-\textbf{value} \\
\midrule
MC Drop.\ vs.\ Deterministic & Dice      & $<0.001$** \\
MC Drop.\ vs.\ Deterministic & AUC       & $<0.001$** \\
MC Drop.\ vs.\ Deterministic & Unc-AUROC & $<0.001$** \\
\addlinespace[2pt]
\midrule
\addlinespace[2pt]
MC Drop.\ vs.\ TTA           & Dice      & $0.142$\phantom{**} \\
MC Drop.\ vs.\ TTA           & Unc-AUROC & $<0.001$** \\
\addlinespace[2pt]
\midrule
\addlinespace[2pt]
MC Drop.\ vs.\ Ensemble      & Dice      & $<0.001$** \\
MC Drop.\ vs.\ Ensemble      & Unc-AUROC & $0.005$** \\
\bottomrule
\end{tabular}
\end{center}
\end{table}

TTA's advantage over MC Dropout is a difference in uncertainty quality,
not segmentation quality. The Dice gap (0.768 vs.\ 0.764) is not
significant, but the Unc-AUROC gap (0.881 vs.\ 0.722) is highly so. This
reinforces the central message: \emph{how} uncertainty ranks errors is
what differentiates these methods, not their raw segmentation outputs.

\subsection{Confidence-aware score distribution}

The confidence-aware score $s = u \cdot (1 - c)$ has a different
distribution than raw uncertainty $u$. Raw uncertainty is right-skewed
(most pixels have low uncertainty, with a long tail). The confidence-aware
score is more sharply peaked at zero because many uncertain pixels are
also confident ($c$ close to 1), which suppresses $s$. The confidence-aware
threshold $\tau_s$ therefore operates in a region where small changes in
$\tau_s$ produce large changes in the deferral set. In practice this
sensitivity is manageable: optimal $\tau_s$ values are stable across
cross-validation folds (CV $<$ 12\%).

\section{Extended calibration analysis}
\label{app:calibration}

\subsection{Temperature scaling optimization}

Temperature scaling minimizes BCE loss on validation logits using L-BFGS
with a maximum of 50 iterations. Convergence typically occurs within 10
iterations. The optimization landscape is convex (BCE loss is convex in
logits, and dividing logits by $T$ is a monotonic reparameterization),
so the global optimum is found reliably.

Learned temperatures: $T = 1.19$ (MC Dropout), $T = 1.35$ (TTA). Both
$T > 1$, confirming mild overconfidence. TTA's higher $T$ reflects the
fact that averaging predictions across augmentations tends to sharpen
the mean prediction, producing slightly more overconfident outputs than
MC Dropout.

\subsection{Bin-level calibration breakdown}

\begin{table}[H]
\centering
\caption{Bin-level calibration analysis for MC Dropout (15-bin ECE). Most
miscalibration concentrates in the low-confidence bins (0.07--0.20), where
the model underestimates accuracy by 1--2 points. Temperature scaling
reduces mid-range gaps by 2--4 percentage points but has minimal effect
on the dominant bins near 0 and 1, which explains why overall ECE is
largely unchanged. Bold marks the largest pre-calibration gap.}
\label{tab:bins}
\scriptsize
\setlength{\tabcolsep}{3.5pt}
\renewcommand{\arraystretch}{1.03}
\begin{center}
\begin{tabular}{l c c c c c}
\toprule
\textbf{Bin range} & \textbf{\% pixels} & \textbf{Avg conf} &
\textbf{Avg acc} & \textbf{Gap} & \textbf{Gap (post-cal)} \\
\midrule
$[0.00, 0.07)$  & 82.3\% & 0.013 & 0.016 & 0.003 & 0.004 \\
$[0.07, 0.13)$  &  2.1\% & 0.095 & 0.112 & 0.017 & 0.015 \\
$[0.13, 0.20)$  &  1.3\% & 0.163 & 0.186 & \textbf{0.023} & 0.019 \\
$[0.20, 0.33)$  &  1.4\% & 0.261 & 0.278 & 0.017 & 0.013 \\
$[0.33, 0.47)$  &  1.1\% & 0.398 & 0.405 & 0.007 & 0.005 \\
$[0.47, 0.60)$  &  0.9\% & 0.531 & 0.522 & 0.009 & 0.008 \\
$[0.60, 0.73)$  &  0.8\% & 0.665 & 0.648 & 0.017 & 0.014 \\
$[0.73, 0.87)$  &  1.0\% & 0.796 & 0.782 & 0.014 & 0.011 \\
$[0.87, 0.93)$  &  1.2\% & 0.898 & 0.891 & 0.007 & 0.006 \\
$[0.93, 1.00]$  &  7.9\% & 0.978 & 0.984 & 0.006 & 0.005 \\
\bottomrule
\end{tabular}
\end{center}
\end{table}

The dominant contribution to ECE comes from the low-confidence bins
(0.07--0.20). The high-confidence bin ($>0.93$) is already well calibrated:
predicted confidence 0.978 matches observed accuracy 0.984. Temperature
scaling reduces mid-range gaps by 2--4 percentage points but has minimal
effect on the dominant low- and high-confidence bins, which contain 90\%
of pixels. This bin-level structure explains why overall ECE moves so
little under calibration and why calibration fails to improve deferral---
the rankings within each bin, which deferral depends on, are untouched.

\section{Comparison with deterministic baseline}
\label{app:deterministic}

The deterministic baseline (single forward pass, no uncertainty estimation)
achieves Dice 0.736, AUC 0.944, and ECE 0.245. The high ECE reflects the
fact that a single sigmoid output, without averaging across multiple
passes, tends to produce extreme probabilities (near 0 or near 1) that
are poorly calibrated.

MC Dropout improves Dice by 2.8 points (0.764 vs.\ 0.736), AUC by 2.1
points, and reduces ECE from 0.245 to 0.038. TTA produces comparable
improvements. The deterministic baseline has Unc-AUROC of exactly 0.500
(random) because, without multiple passes, no meaningful uncertainty
estimate exists. This confirms that uncertainty estimation requires
stochasticity at test time---whether from dropout masks, input transforms,
or separate models.

\section{Notation reference}
\label{app:notation}

\begin{table}[H]
\centering
\caption{Summary of notation used throughout the paper.}
\label{tab:notation}
\scriptsize
\setlength{\tabcolsep}{3.5pt}
\renewcommand{\arraystretch}{1.05}
\begin{center}
\begin{tabular}{l l}
\toprule
\textbf{Symbol} & \textbf{Meaning} \\
\midrule
$\mathbf{x}$       & Input image ($\mathbb{R}^{H \times W \times C}$) \\
$\mathbf{y}$       & Ground truth mask ($\{0,1\}^{H \times W}$) \\
$\hat{\mathbf{p}}$ & Predicted probability map ($[0,1]^{H \times W}$) \\
$\hat{y}_{ij}$     & Hard prediction: $\mathbb{1}[\hat{p}_{ij} > 0.5]$ \\
$\mathbf{u}$       & Uncertainty map ($\mathbb{R}_{\geq 0}^{H \times W}$) \\
$T$                & Number of MC Dropout forward passes \\
$K$                & Number of TTA augmentations \\
$\theta_t$         & Network weights under dropout mask $t$ \\
$\tau_k$           & $k$-th geometric augmentation \\
$\bar{p}_{ij}$     & Mean prediction across passes \\
$\mathrm{MI}_{ij}$ & Mutual information (epistemic uncertainty) \\
$H(p)$             & Binary entropy, $-p\log p - (1-p)\log(1-p)$ \\
$\delta_{ij}$      & Deferral decision (1 = accept, 0 = defer) \\
$\tau$             & Deferral threshold (global) \\
$\tau_n$           & Deferral threshold for image $n$ (adaptive) \\
$c_{ij}$           & Prediction confidence: $2|\hat{p}_{ij} - 0.5|$ \\
$s_{ij}$           & Confidence-aware score: $u_{ij} \cdot (1 - c_{ij})$ \\
$T_{\mathrm{cal}}$ & Temperature scaling parameter \\
\bottomrule
\end{tabular}
\end{center}
\end{table}
